% Transcription — fonds Alexandre Grothendieck, shelfmark 156-7, batch 2 (pages 21-40).
% Established from the facsimile archives/batches/156-7/batch-02.pdf (a copy of
% 156-7.pdf fetched through the project's relay; no cover sheet in the batch,
% so PDF page k = folder page 20+k), per the skill transcribe-grothendieck. The
% folder is 113 pages in six batches, transcribed in parallel; this is the
% second. Batch 1 of this folder did not exist when this batch was drafted, so
% the axioms and lemmas it states (At C, At D, At D_0, At CL (spéc), At L 1-3,
% the lemma of page 15, the NB of page 17, At supp on page 18) are cited here
% only as this batch cites them. The notation of the second draft (folder
% 156-6, GF VI) is kept: 𝔉 \mathfrak{F}, 𝓜 \mathcal{M}, ℒ \mathcal{L},
% compatibility (crossed chevrons) \between, disjointness \parallel, ≪ \ll,
% interior refinement \overset{\circ}{\ll}, « omb », « Omb » in roman.
%
% Pass: Opus 5.5 (claude-opus-5-5), 2026-09-26 — first pass, unchecked against the pages by a human.
%
% Pages skipped: none. Pages summarised: none. All twenty leaves are his
% mathematics, in ink, fast drafting with many strikes and slanted margin
% notes (read only in part on pages 24, 34, 39).
%
% His pagination 21-40 stands at the head of each leaf and coincides with the
% archivists' (noted once, on page 21). No date in his hand on these twenty
% pages; the 23-26/06/1986 of the catalogue title is not repeated here.
%
% Numbered runs, recorded where they stand. Page 21: Lemme 1 (x_P = ∅ iff
% x ∉ omb(P)) and Lemme 2 (F_P = ∅ iff F ∥ P), the condition (AL) and its
% consequence (A). Page 22: At CL (pol); the implications a)-e) among At L 2,
% At L 3, At CL, At C, At D, At D_0 (margin: « Voir diagr. refait p. 45 »).
% Page 23: the « Récapitulation des axiomes »: At 1-At 7 (« gros œuvre »), At C,
% At D, At D_0, At C (spéc), At C (pol), At Fil., At supp., At L 1-3, At CL
% (spéc), At CL (pol). Page 24: formulas (1)-(4). Pages 25-33, a new run in
% terms of (𝓜, ≤, ≪, ∥ on ℒ): ML 0, ML 1 (p. 25; ML 2, ML 3 struck there and
% restated p. 26), a new Lemme 1 (p. 25), Cor, ML 2, ML 3 (p. 26), Lemme 2
% (p. 27), Corollaires 1-3 (p. 28), Corollaire 4 (p. 29), a « ML 0 » for the
% relation ∥ (p. 29, not the ML 0 of p. 25), formula (1) (p. 29), Lemme 3
% (p. 29), ML 4, Lemmes 4-5 (p. 30), Lemme 6 with a Cor, Lemme 7 (p. 32),
% Lemme 8, a struck « Lemme 9 », the Théorème-Scholie (p. 33). Page 34
% renumbers the axioms ML 0-ML 5 (the old ML 0 for ∥ becomes ML 4, the old
% ML 4 becomes ML 5); page 36 speaks of « quatre axiomes (cinq avec (ML 0))
% ML 1 - ML 4 »; ML 6 and ML 0 pol on page 37. Pages 39-40, restart with a
% dense part Λ ⊂ 𝓜: MΛ 0-MΛ 3, Lemmes A, B, C and a Corollaire (restating
% Lemme 1 p. 25 and Cors. 1-3 p. 28, as his margin says), MΛ 0', and a struck
% MΛ 4. His names for the atelier axioms At_𝓜 4 a), At_𝓜 4 b), At_𝓜 5, At_𝓜 6
% (pp. 26-32; « page 6 » cited for At_𝓜 4 a)) are kept as written. Internal
% cross-references, kept as plain page numbers: p. 45 (p. 22), p. 17 (p. 22),
% p. 15 (p. 21), p. 6 (p. 26), p. 29 (pp. 32, 35, 37), p. 13 (p. 37), p. 18 and
% « GF VI, p. 52 » (p. 38; the second draft's page on the ensemblist atelier
% and At ens 1), p. 25 and p. 28 (p. 40).
%
% What the batch is about. The third draft of Analysis situs sorts the
% compatibility-disjunction axioms of an atelier and then asks how much of an
% atelier is recovered from its multistrates. Pages 21-23 finish the proof
% that At CL (spéc) gives the compatibility criterion by refinements with
% X'_L = Y'_L = ∅, prove two lemmas on x_P and F_P, introduce At CL (pol),
% chart the implications among At L 1-3, At CL, At C, At D, At D_0, and
% recapitulate all the axioms. Pages 24-38: given (𝓜, ≤, ≪) and a symmetric
% antireflexive relation ∥ on the ≪-minimal elements ℒ, compatibility on 𝓜 is
% defined by (1) p. 29 (X ≬ Y iff all lieux of X and of Y outside omb(X ∩ Y)
% are disjoint); with omb(X)° the interior shadow, a chain of lemmas (strata
% detected by interior lieux, transitivity of interior refinement, heredity
% of compatibility under ≤) leads to the Théorème-Scholie (pp. 33-35): the
% data come from a « spécieux » atelier satisfying At 1-7, At L 1-3 and At CL
% (spéc) iff ML 0-ML 5 hold; that atelier is the set Φ of closed parts of 𝓜
% made of pairwise compatible elements, and every such atelier is a sub-atelier
% of Φ; a polyhedral variant (ML 0 pol), the freedom in the choice of ∥
% (any graph Γ ⊃ Γ_0 on ℒ, « Ça fait de la marge ! »), and the question of
% ensemblist ateliers (x ∥ y iff x ≠ y is not enough). Pages 39-40 start over
% with a dense part Λ in place of ℒ and ≬ given on 𝓜.
%
% Reading hazards fixed once here. His underlined capital Φ (p. 35) is
% \underline{\Phi}. Where he draws ≤ or ≪ as vertical or oblique strokes
% between letters, the configuration is set as a tikzcd with headless
% labelled edges (pp. 27, 29-31, 34), as in 156-6 batch 4. Variances kept and
% flagged in notes: X ∥ Y for X' ∥ Y' in (A) (p. 21); ML 0 cited for ML 1
% (p. 29); Y' for Y (p. 26); {{X}} for X̃ (p. 36); « ML 4 est vide » where
% ML 5 is meant (p. 34); « Z ≪ X' » kept in X̃^Y (p. 40).
%
% Apparatus: 138 \ill{} and 20 \uncertain{}; hardest pages 23-24, 34, 38-40.
\documentclass[11pt,a4paper]{article}
\input{../preamble/grothendieck}

\folder{156-7}
\batch{2}
\pages{21}{40}
\dating{1986}
\watermark{Édition de démonstration}
\foldertitle{[Chapitre] VII. Analysis situs (troisième mouture) : notes manuscrites (23-26/06/1986)}

\begin{document}

\page{21}
\note{son numéro « 21 » en haut de la page ; sur tout ce lot sa pagination, en tête de chaque feuillet, coïncide avec celle des archivistes, comme dans les deux premières moutures}

\ill{} \ill{} $x \in \mathrm{omb}(X') \subset \mathrm{omb}(X)$, $y \in \mathrm{omb}(Y') \subset \mathrm{omb}(Y)$, \ill{} $x \parallel y$. Or la condition $X'_L = \emptyset$ implique $x_L = \emptyset$ (car $x_L \ll X'_L$), et \struck{\ill{}} cela signifie ici que $x \notin \mathrm{omb}(L)$ (cf lemme) et itou pour $y$. Donc $x \in \mathrm{omb}(X) \smallsetminus \mathrm{omb}(L)$, $y \in \mathrm{omb}(Y) \smallsetminus \mathrm{omb}(L)$, et si on a At CL (spéc), on conclut bien $x \parallel y$,
\note{la page commence au milieu d'un raisonnement engagé à la page 20 (lot 1) : la démonstration que At CL (spéc) entraîne le critère de compatibilité en termes de raffinements $X' \ll X$, $Y' \ll Y$ avec $X'_L = Y'_L = \emptyset$}

\struck{\ill{}} Donc la condition
\[
  (\mathrm{AL}) \qquad \forall\, x \in \mathrm{omb}(X) \smallsetminus \mathrm{omb}(L),\ y \in \mathrm{omb}(Y) \smallsetminus \mathrm{omb}(L), \ \text{on a}\ x \parallel y
\]
\note{dans le premier $\mathrm{omb}(L)$, le $L$ est écrit en surcharge sur une autre lettre}
implique
\[
  (\mathrm{A}) \qquad \forall\, X' \in \mathrm{Omb}(X),\ Y' \in \mathrm{Omb}(Y),\ X'_L = Y'_L = \emptyset, \ \text{on a}\ X \parallel Y
\]
\note{« $X \parallel Y$ » est sur la page ; on attend $X' \parallel Y'$}
et l'inverse est trivial, comme on voit en prenant $X, Y \in \mathcal{L}$. On a utilisé ici le
\note{« $X, Y \in \mathcal{L}$ » est sur la page ; on attend $X', Y' \in \mathcal{L}$}

\emph{Lemme 1} Soient \struck{$L \leq F$} $P \leq Q$ (en l'occurrence \struck{\ill{}} $L \leq X$ ou $L \leq Y$) $x \in \mathrm{omb}(Q)$, alors \struck{$x \parallel Q$}
\[
  x_P = \emptyset \iff x \notin \mathrm{omb}(P) \ \text{i.e.}\ x \not\ll P
\]
(\struck{\ill{}})
\note{dans $\mathrm{omb}(P)$ et dans $x \not\ll P$, le $P$ est écrit en surcharge sur une autre lettre ; la parenthèse finale ne contient qu'une formule biffée, non lue}

En effet, \struck{$x \ll P \iff x_P \neq$} pour la première équivalence, passant aux négations, on \struck{trouve} \add{doit} prouver
\[
  x \ll P \iff x_P \neq \emptyset_{\mathfrak{F}} ,
\]
et comme $x_P \ll x$ et $x$ minimal dans $\mathcal{M} \smallsetminus \lbrace \emptyset_{\mathfrak{F}} \rbrace$ pour $\ll$, $x_P \neq \emptyset_{\mathfrak{F}}$ signifie aussi $x_P = x$ i.e.\ $x \ll P$. OK.

\struck{D'autre part, $x \parallel P \Rightarrow x \not\ll P$ est trivial, \ill{}} Je profite de l'occasion pour expliciter \struck{la définition} (moyennant At D)

\emph{Lemme 2} Soient $P \leq Q$, et $F \ll Q$. Conditions équivalentes
\[
  F_P = \emptyset \iff F \parallel P .
\]
\note{l'énoncé du lemme 2 est marqué d'un trait vertical dans la marge}
On utilise que At D (ou plutôt, la prop.\ qui suit, \ill{} p.~15) à $P, Q$ \quad $P \leq P$, $F \ll Q$), pour avoir $\Longrightarrow$ ; l'autre est vrai avec les seuls axiomes At 1 -- At 7.
\note{la page 15 est au lot 1}

\page{22}
On introduit donc
\note{« On introduit donc » est suivi d'un signe = ; l'axiome qui suit est souligné}

\emph{At CL (pol)} (Critère de compatibilité par les lieux, pour ateliers polyédraux)

Notons
\note{la première ligne du tableau, biffée, se lit : At D$_0$ + At L 1 + At CL (spéc) $\Longleftrightarrow$ At L 2 $\Longrightarrow$ At D, avec à droite une double flèche descendante vers « At D$_0$ », barrée elle aussi}
\[
  \begin{array}{ll}
    a) & \text{At L 2} + \text{At CL (spéc) (ou pol)} \ \Longleftrightarrow\ \text{At L 2} + \text{At C (spéc) (ou pol)} \\
    b) & \text{At D}_0 + \text{At L 1} + \text{At CL (spéc) (ou pol)} \ \Longrightarrow\ \text{At L 2} \\
    c) & \text{At L 2} + \text{At L 3} \ \Longrightarrow\ \text{At D} \ \Longrightarrow\ \text{At D}_0 \\
    d) & \text{At L 3} + \text{At CL (spéc)} \ \Longrightarrow\ \text{At C} \\
    e) & \text{At D} \ \Longrightarrow\ \text{At L 3}
  \end{array}
\]
\note{les cinq lignes a)--e) sont réunies par une accolade à gauche. Dans a), un « At » est biffé en tête ; dans b), « At D$_0$ + » est écrit au-dessus d'un mot biffé, et un « At D » est biffé après « At L 2 », d'où part une double flèche descendante vers le « At D$_0$ » de c) ; dans c), « At D » est récrit sur un premier « At D » biffé, et une double flèche remonte de « At C » (ligne d)) vers « At D » ; à droite de d), un groupe de mots est biffé ; dans e), le « At » de « At L 3 » est en surcharge ; sous e), « At L 1 + A » biffé}

\marginal{Voir \uncertain{diagr.}\ refait p.~45}
\note{note écrite en oblique à droite des lignes c)--e) ; la page 45 est au lot 3}

Pratiquement, je retiens ceci, concernant l'ensemble des axiomes de compatibilité-disjonction [et de divisibilité]
\[
  \begin{array}{lll}
    C \Rightarrow D \Rightarrow D_0 , & \text{At C (spéc)} & \text{At CL (spéc)} \\
    \quad \Downarrow & \text{At C (pol)} & \text{At CL (pol)} \\
    {[L 1]}\ L 2,\ L 3 & &
  \end{array}
\]
\note{la double flèche descend de « D » vers « L 2, L 3 »}

1) Du point de vue des axiomes inconditionnels (sans \add{restriction} spéc ou pol \ldots), les deux axiomes \ill{} L 2, L 3 de disjonction impliquant les lieux impliquent D donc D$_0$, mais non C (\struck{\ill{}} \ill{} de compatibilité, \ill{} \ill{} n'intervenant pas : une disjonction). Pour déduire C d'axiomes pertinents sur les lieux, il nous faut des At CL spéc ou At CL (pol). À ce sujet, on note
\[
  \text{At L 2} + \text{At L 3} + \text{At CL (spéc)} \quad
  \begin{array}{l}
    \text{implique tous les axiomes de compatibilité-} \\
    \text{disjonction (style spéc) :} \\
    \text{At C (spéc) \quad par (a)} \\
    \text{At D \quad par (c) d'où D}_0 \\
    \text{At C \quad par NB p.~17 ou par d)}
  \end{array}
\]
\note{la page 17 est au lot 1}

\struck{On peut \ill{} \ill{} \ill{} At L} « pol » ald « spéc »
\note{« ald » : au lieu de ; ces mots, en bas à gauche de la page, sont écrits sous une ligne biffée}

\page{23}
L'axiome de divisibilité At L 1 est à part, il n'intervient pas encore dans ce qui précède, mais on a
\[
  \begin{array}{l}
    \text{At L 1} + \text{At L 2} + \text{At L 3} + \text{At CL (spéc)} \\
    \quad \Updownarrow \\
    \text{At D}_0 + \text{At L 1} + \text{At L 3} + \text{At CL (spéc)} \\
    \quad \Downarrow \\
    \text{At D} + \text{At L 1} + \text{At CL (spéc)} \qquad \text{At C (spéc)}
  \end{array}
\]
implique tous les axiomes de compatibilité --- en \uncertain{supposant}, i.e.\ style « spéc » --- plus bien sûr At L 1 lui-même.
\note{cette dernière phrase est écrite à droite des deux premières lignes du tableau}
\note{sous la première ligne, « At D$_0$ » biffé. Sur la dernière ligne, « At C (spéc) » est écrit au-dessus, à droite, comme une variante ; au-dessous vient une ligne mal lue : « et \ill{} At D + At L 1 + At L 2 + \ill{} \ill{} « pol » ald « spéc » », le « At D » de tête étant repassé. Une note oblique de la marge gauche, d'où partent deux flèches vers ces deux lignes, se lit : « \struck{\ill{}} \ill{} deux axiomes relatifs At D, At C (sp) et deux axiomes relatifs à L »}

Récapitulation des axiomes (à l'exclusion des \struck{\ill{}} axiomes d'extension d'une subdivision, version générale, ou « spéc », ou « pol »\ldots)

\begin{itemize}
  \item[At 1] Sup des familles filtrantes majorées
  \item[At 2] Critère de $\leq$ par strates
  \item[At 3] Critère de \struck{disjonction} compatibilité par strates
  \item[At 4] Critère de $\ll$ par strates
  \item[At 5] Critère des raffinements entre sous-figures (\ill{})
  \item[At 6] Axiome de la plus petite strate, ou des raffinements intérieurs de strates $X \mathrel{\overset{\circ}{\ll}} Y$
  \item[At 7] Transitivité des raffinements intérieurs
\end{itemize}
\marginal{axiomes du « gros œuvre »}
\marginal{$\leq$}
\marginal{\uncertain{compatibilité} \ill{} de $\leq$ et $\ll$}
\note{accolades et mots de la marge gauche : « $\leq$ » réunit At 1 -- At 3, une seconde accolade At 4 -- At 7, une grande accolade, dans la marge, At 1 -- At 7 avec « axiomes du “gros œuvre” » écrit de bas en haut. La parenthèse de At 5 contient une formule entre sous-figures, non lue avec assez de sûreté, du type « $G \leq F$, $G \leq H$, $F \ll G \Leftrightarrow F \ll G$ »}

\begin{itemize}
  \item[At C] Critère de compatibilité de raffinements de sous-figures.
  \item[At D] Critère de disjonction de --- --- ---
  \item[At D$_0$] \struck{disjonction} stabilité de $\parallel$ par $\ll$
  \item[At C (spéc)] Critère de compatibilité dans les ateliers spécieux
  \item[At C (pol)] --- --- --- polyédraux
\end{itemize}
\marginal{axiomes de compatibilité, disjonction}
\note{des flèches descendantes vont de At C à At D et de At D à At D$_0$}

\begin{itemize}
  \item[At Fil.] Existence des Sup filtrants de \add{une famille loc.\ finie de} raffinements
  \item[At supp.] \ldots
\end{itemize}
\note{ces deux lignes sont entourées d'un cadre, d'où part une longue flèche qui descend le long de la marge droite, passe devant At CL et sort au bas de la page}

\begin{itemize}
  \item[At L 1] Axiome de divisibilité
  \item[At L 2] Critère de disjonction par les lieux
  \item[At L 3] Axiome de disjonction pour les lieux
  \item[At CL (spéc)] Critère de compatibilité par les lieux (at.\ spécieux)
  \item[At CL (pol)] --- --- --- (--- polyédraux)
\end{itemize}
\marginal{axiomes concernant les lieux}

\page{24}
Je \uncertain{reviens} sur la description d'un atelier spécieux ou polyédral, satisfaisant L 1, L 2, L 3 plus \struck{At L (spéc)} ou At CL (spéc) (ou, au choix, At CL (pol)), en termes de $(\mathcal{M}, \leq, \ll, \between)$\,*.

\marginal{* \ill{} \ill{} \ill{} au moins \ill{} \struck{\ill{}} \ill{} \ill{} \ill{} de lieux\ldots}
\note{note écrite en oblique dans la marge gauche, appelée par l'astérisque qui suit $(\mathcal{M}, \leq, \ll, \between)$ ; lue seulement en partie}

Le point ici, c'est que le critère At CL (spéc) \ill{} réduit, pour connaître $\underset{\mathcal{M}}{\between}$, à la connaissance de $\underset{\mathcal{L}}{\parallel}$, i.e.\ à celle d'une structure de graphe sur $\mathcal{L}$ (NB dans le cas d'un atelier ensembliste, je rappelle que $x \parallel y \Leftrightarrow x \neq y$). Donc, la question est, partant de \struck{$(\mathcal{M}, \leq, \ll)$}, \struck{d'où $\mathcal{L}$}
\note{devant $\underset{\mathcal{L}}{\parallel}$, un signe de compatibilité est biffé}
\[
  (1) \qquad (\mathcal{M}, \leq, \ll)
\]
d'où
\[
  (2) \qquad \mathcal{L} = \lbrace x \in \mathcal{M} \mid x \ \text{minimal pour}\ \ll \rbrace ,
\]
et d'une relation symétrique antiréflexive $x \parallel y$ ($x, y \in \mathcal{L}$) (\ill{} $x \neq y$) sur $\mathcal{L}$, quelles conditions \ill{} pour que ça corresponde bien à un \ill{} spécieux, satisfaisant \struck{At} (ou plus At CL (spéc)) les conditions L 1, L 2, L 3.

\marginal{\ill{}}
\note{une seconde note, longue, écrite en oblique sur sept ou huit lignes dans la marge gauche en regard de (2)--(3), n'a pas été lue ; on y distingue « figures », « lieux », « compatibles », « atelier », « L (?) » et « $\mathcal{M}_{\leq}$ »}

Il est entendu qu'on va \emph{définir} $\underset{\mathcal{M}}{\between}$ par
\[
  (3) \qquad X \between Y \iff \forall\, x \in \mathrm{omb}(X) \smallsetminus \mathrm{omb}(L),\ y \in \mathrm{omb}(Y) \smallsetminus \mathrm{omb}(L), \ \text{on a}\ x \parallel y ,
\]
où
\[
  L = X \cap Y = \operatorname*{Sup}_{Z \in \widetilde{X} \cap \widetilde{Y}} Z \qquad \Bigl(\text{donc}\ \mathrm{omb}(L) \overset{\mathrm{def}}{=} \bigcup_{Z \in \widetilde{X} \cap \widetilde{Y}} \mathrm{omb}(Z)\Bigr) .
\]
Pour commencer, il faut \struck{vérifier} \add{écrire} les conditions préliminaires
\[
  (4) \qquad X \between Y,\ X' \leq X,\ Y' \leq Y \Longrightarrow X' \between Y'
\]
\struck{et nous devrions le faire sous des formes plus}

\page{25}
D'abord, \struck{précisons} dégageons quelques conditions concernant $\leq$, $\ll$, \struck{\ill{}} indépendamment de $\parallel$.
\[
  (\mathrm{ML}\ 0) \qquad X \leq Y \Longrightarrow X \ll Y \qquad \text{pour mémoire}
\]
\note{le « 0 » de ML 0 est repassé}

(ML 1) Soit $X' \leq X \geq X''$, soit $x \in \mathcal{L}$ avec $x \ll X'$, $x \ll X''$, alors $\exists\, Z \in$ \ill{} avec $x \ll Z$ et $Z \leq X', X''$.
\note{sur la page, $X''$ est écrit sous $X$, relié à lui par un signe $\geq$ vertical}
\marginal{[on doit pouvoir supposer $x \mathrel{\overset{\circ}{\ll}} Z$ \ill{} $Z \in \mathcal{M}$ \uncertain{quelc.}\,?]}
\note{note entre crochets à droite de ML 1, rattachée par une flèche à « $x \in \mathcal{L}$ »}

[I.e.\ \struck{$X' \cap X''$} la « figure intersection » $X' \cap X'' = \widetilde{X}' \cap \widetilde{X}''$ est un Inf de $X', X''$ pour $\ll$, \ill{} \ill{} \ill{} $\mathcal{L}$)

\marginal{Il faudra renforcer, avec $X' \ll X$, $X'' \leq X$, \ill{} les propriétés \ill{} Inf$(X', X'')$}
\note{note oblique dans la marge gauche, cernée d'un trait qui la rattache à ML 1, et marquée d'un double trait vertical}

\struck{(\ill{})}
Posons \struck{pour $\mathrm{omb}(X) = \mathrm{omb}$} \add{$X \in \mathcal{M}$}
\[
  \begin{aligned}
    \mathrm{omb}(X) &= \lbrace x \in \mathcal{L} \mid x \ll X \rbrace \\
    \mathrm{omb}(X)^{\circ} &= \mathrm{omb}(X) \smallsetminus \bigcup_{Y < X} \mathrm{omb}(Y)
  \end{aligned}
\]
Posons
\[
  x \mathrel{\overset{\circ}{\ll}} X \overset{\mathrm{def}}{\iff} x \in \mathrm{omb}(X)^{\circ} \quad \text{i.e.} \quad x \ll X, \ \text{et}\ x \ll X' \leq X \Rightarrow X' = X .
\]
On définit de même
\[
  Y \mathrel{\overset{\circ}{\ll}} X \iff Y \ll X, \ \text{et}\ Y \ll X' \leq X \Rightarrow X' = X ,
\]
i.e.\ $X$ minimal dans $\widetilde{X}^{Y} = \lbrace X' \in \widetilde{X} \mid Y \ll X' \rbrace$.

\struck{(ML 2) Soit $X \in \mathcal{M}$. Alors $\mathrm{omb}(X)^{\circ} \neq \emptyset$}

\struck{(ML 3) Soit $X \in \mathcal{M}$, $x \in \mathcal{L}$, $x \ll X$ \ill{} $\exists\, Z$, \ill{} $\leq X$.}
\note{ces deux énoncés sont barrés de traits obliques ; ML 2 et ML 3 sont repris à la page 26}

\emph{Lemme 1} Soient \struck{$X, Y \in \mathcal{M}$} \ill{} Conditions équivalentes.
\marginal{\struck{At L 1} (ML 2) \uncertain{avant} le lemme 1}
\note{note de la marge gauche, reliée par une accolade à ML 2 -- ML 3 et au lemme 1}
\begin{itemize}
  \item[(i)] $X \leq Y$
  \item[(ii)] $X \ll Y$
  \item[(ii bis)] $\mathrm{omb}\, X \subset \mathrm{omb}\, Y$
  \item[(iii)] $\mathrm{omb}(X)^{\circ} \subset \mathrm{omb}(Y)$
  \item[(iv)] $\mathrm{omb}(X)^{\circ} \cap \mathrm{omb}(Y) \neq \emptyset$
\end{itemize}
\note{à droite, « $\Downarrow$ via (ML 2) » entre (iii) et (iv), sous un premier renvoi biffé « (ML 2) »}

Il suffit de prouver (iv) $\Longrightarrow$ (i). Or si $x \in \mathrm{omb}(X)^{\circ} \cap \mathrm{omb}(Y)$, $\exists$ par ML 1 $\exists\, Z$, $Z \leq X, Y$, avec $x \ll Z$. On a $x \ll Z \leq X$, d'où $Z = X$ puisque $x \mathrel{\overset{\circ}{\ll}} X$, donc $X \leq Y$, cqfd.

\page{26}
\emph{Cor} Si $\mathrm{omb}(X)^{\circ} \cap \mathrm{omb}(Y)^{\circ} \neq \emptyset$, on a $X = Y$ [i.e.\ $X \neq Y \Rightarrow \mathrm{omb}(X)^{\circ} \cap \mathrm{omb}(Y)^{\circ} = \emptyset$]

\emph{ML 2} \quad $\forall\, X \in \mathcal{M}$, $\mathrm{omb}(X)^{\circ} \neq \emptyset$. i.e.\ $\exists\, x \in \mathcal{L}$, $x \mathrel{\overset{\circ}{\ll}} X$

\emph{ML 3} \quad $\forall\, X \in \mathcal{M}$, les $\mathrm{omb}(Y)^{\circ}$ ($Y \in \widetilde{X}$) recouvrent $\mathrm{omb}(X)$, i.e.\ \struck{$\forall$} $\forall\, x \in \mathcal{L}$, $x \ll X$, $\exists$ \struck{$X$} $Y \leq X$, avec $x \mathrel{\overset{\circ}{\ll}} Y$.
\marginal{At L 1}
\marginal{At$_{\mathcal{M}}$ 4 a)}
\marginal{\ill{} ML 2 et ML 3 \uncertain{signifient} que les $\mathrm{omb}(Y)^{\circ}$ forment une partition de $X$.}
\note{ML 2 et ML 3 sont marqués d'un trait vertical ; « At L 1 » est écrit dans la marge en regard de ML 2, « At$_{\mathcal{M}}$ 4 a) » en regard de ML 3, et la troisième note, oblique, au-dessous}

\emph{NB} Par le cor.\ ci-dessus, cet $Y$ est unique. Par le lemme, si \struck{$x \mathrel{\overset{\circ}{\ll}}$} $Y' \leq X$, et $x \ll Y'$, comme $x \in \mathrm{omb}(Y)^{\circ} \cap \mathrm{omb}(Y')$ \ill{} on a $Y \leq Y'$ par le lemme, donc $Y'$ est le plus petit élément de $\widetilde{X}^{x}$. Plus généralement :
\note{on attend : $Y$ est le plus petit élément de $\widetilde{X}^{x}$ ; « $Y'$ » est sur la page}

\struck{Lemme 2 Soit $X, Y \in \mathcal{M}$, $Y \ll X$. Alors}

On n'a pas jusqu'à présent utilisé la définition des lieux $\mathcal{L}$ de $\mathcal{M}$, comme formé d'él.\ minimaux pour $\ll$. Dans \uncertain{tout} modèle avec un $\mathcal{L}$ quelconque, pourrions \ill{} supposer ML 2, ML 3 pour cet $\mathcal{L}$. Si on prenait $\mathcal{L} = \mathcal{M}$, ML 2 est trivial, ML 3 signifie que $\forall$ \struck{\ill{}} $X \in \mathcal{M}$ et $Z \ll X$, $\exists\, Y \in \widetilde{X}$ tel que $Z \mathrel{\overset{\circ}{\ll}} Y$, i.e.\ At$_{\mathcal{M}}$ 4 a) (page 6). Mais il faudrait dans ML 1 \struck{\ill{}} \add{généraliser}, en remplaçant $x$ par un $Z$ \add{$\in \mathcal{M}$} quelconque.
\note{la page 6 est au lot 1}

Mais il devrait être possible de le prouver, via ML 3 \add{et ML 2} \struck{pour} $Z$. L'idée est \struck{\ill{}} \add{(ML 2)} $Z \ll X$, \add{de} prendre $x \in \mathcal{L}$, $x \mathrel{\overset{\circ}{\ll}} Z$, et l'unique $Y \in \widetilde{X}$ tel que $x \mathrel{\overset{\circ}{\ll}} Y$ (ML 3), et il \struck{par} faudrait prouver que $Z \mathrel{\overset{\circ}{\ll}} Y$. Ceci se \uncertain{décompose} en deux : si $x \ll Y \leq X$, on a $Z \ll Y$ et donc
\[
  x \mathrel{\overset{\circ}{\ll}} Y \ \text{ssi}\ Z \mathrel{\overset{\circ}{\ll}} Y .
\]
\note{dans « $Z \mathrel{\overset{\circ}{\ll}} Y$ » de la ligne précédente, le $Z$ et le $Y$ sont écrits en surcharge ; dans « $x \ll Y \leq X$, on a $Z \ll Y$ », le premier $\ll$ est doublé}

\page{27}
Donc on \ill{} prouver

\emph{Lemme 2} Soit $Z \ll X \geq X'$, soit $x \in \mathcal{L}$, tel que $x \ll Z$, $x \ll X'$. Si $x \mathrel{\overset{\circ}{\ll}} Z$, on a $Z \ll X'$, et donc
\[
  x \mathrel{\overset{\circ}{\ll}} X' \iff Z \mathrel{\overset{\circ}{\ll}} X' .
\]
\note{sur la page, $X'$ est écrit sous $X$, relié par un signe $\geq$ vertical, et repassé ; un petit « $x$ » est écrit au-dessus du premier $\ll$ ; le $\ll$ de « $Z \ll X'$ » est surchargé}

\emph{On admettra} ML 1 \emph{sous forme généralisée}, \ill{} ici que $\exists\, Z' \leq Z$, \struck{$Z \ll X'$} $Z' \ll X'$ et $x \ll Z'$ :

\begin{tikzcd}[column sep=small, row sep=small]
 & Z \arrow[r, no head, "\ll" description] & X \\
 & Z' \arrow[u, no head, "\leq" description] \arrow[r, no head, "\ll" description] & X' \arrow[u, no head, "\leq" description] \\
x \arrow[ur, no head, "\ll" description] & &
\end{tikzcd}

Ceci dit, si $x \mathrel{\overset{\circ}{\ll}} Z$, on a $Z' = Z$, donc $Z \ll X'$. Si de plus $x \mathrel{\overset{\circ}{\ll}} X'$, alors \emph{a fortiori} $Z \mathrel{\overset{\circ}{\ll}} X'$ (tautologique), \struck{$Z' \mathrel{\overset{\circ}{\ll}} x \ll$} prouvons l'inverse \struck{$Z$}
\[
  x \mathrel{\overset{\circ}{\ll}} Z \mathrel{\overset{\circ}{\ll}} X' \Longrightarrow x \mathrel{\overset{\circ}{\ll}} X' .
\]
Soit en effet \struck{$X'' \leq$} $X''$ avec $x \ll X'' \leq X'$, par ce qui précède, appliqué à $X'' \leq X$ (car $X' \leq X$) on a $Z \ll X'' \leq X'$, et comme $Z \mathrel{\overset{\circ}{\ll}} X'$ on conclut $X'' = X'$, cqfd.

Ceci prouve donc que $\forall\, Z$ tel que $Z \ll X$, $\exists\, X' \leq X$ avec $Z \mathrel{\overset{\circ}{\ll}} X'$ (en utilisant ML 1 renforcé, ML 2 et ML 3) (c'est At$_{\mathcal{M}}$ 4 a)). On trouve en \ill{} temps

\page{28}
\struck{\emph{Corollaire} Soit $Z \ll Y \ll X$ $X \ll Y \ll Z$. Conditions équivalentes (i) $Z \mathrel{\overset{\circ}{\ll}} Y$ et $Z \ll X$}
\note{début de corollaire annulé par un grand trait courbe et des ratures}

\emph{Corollaire 1} Pour $Z \ll X$, $\exists\, X'$ avec $Z \mathrel{\overset{\circ}{\ll}} X' \leq X$.

\emph{Corollaire 2} Soit $Z \ll X \geq X'$. Conditions équivalentes
\note{$X'$ est écrit sous $X$, relié par un $\geq$ vertical. Deux premières conditions sont biffées : « (i) $Z \mathrel{\overset{\circ}{\ll}} X$ » et « (ii) $\mathrm{omb}(Z)^{\circ} \subset \mathrm{omb}(X)$ »}
\begin{itemize}
  \item[(i)] $Z \ll X'$
  \item[(i bis)] $\mathrm{omb}(Z) \subset \mathrm{omb}(X')$
  \item[(ii)] $\mathrm{omb}(Z)^{\circ} \subset \mathrm{omb}(X')$
  \item[(iii)] $\mathrm{omb}(Z)^{\circ} \cap \mathrm{omb}(X') \neq \emptyset$
\end{itemize}
\note{« $\Downarrow$ par ML 2 » à droite, entre (ii) et (iii)}

Il suffit de voir (iii) $\Longrightarrow$ (i), et c'est ce qu'énonce le lemme précédent.

\emph{Corollaire 3} \struck{\ill{}} Les notations précédentes, $Z \ll X$, conditions équivalentes
\begin{itemize}
  \item[(i)] $Z \mathrel{\overset{\circ}{\ll}} X'$
  \item[(ii)] $\mathrm{omb}(Z)^{\circ} \subset \mathrm{omb}(X')^{\circ}$
  \item[(iii)] $\mathrm{omb}(Z)^{\circ} \cap \mathrm{omb}(X')^{\circ} \neq \emptyset$
  \item[(iv)] $X'$ est minimal dans $\widetilde{X}^{Z} = \lbrace X' \in \widetilde{X} \mid Z \ll X' \rbrace$
  \item[(v)] $X'$ est un plus petit élément dans $\widetilde{X}^{Z}$
\end{itemize}
\note{le « (v) » est écrit en surcharge sur un « (iv) » ; devant, une ligne biffée « (i) $\Longrightarrow$ (iii) \ill{} »}

On a
\[
  \begin{array}{c}
    (v) \\ \Downarrow \\ (i) \Longleftrightarrow (iv)
  \end{array}
  \qquad (ii) \Longrightarrow (iii) \qquad \text{triviales.}
\]
On va prouver (iii) $\Longrightarrow$ (v) et (i) $\Longrightarrow$ (ii), ce qui suffira.

(i) $\Longrightarrow$ (ii) a été vu : si $x \in \mathrm{omb}(Z)^{\circ}$ (existe par ML 2) \ill{} $x \mathrel{\overset{\circ}{\ll}} Z$, et $Z \mathrel{\overset{\circ}{\ll}} X'$, on a (lemme 2) $x \mathrel{\overset{\circ}{\ll}} X'$ i.e.\ $x \in \mathrm{omb}(X')^{\circ}$, donc (ii).

(iii) $\Longrightarrow$ (v) Soit $x \in \mathrm{omb}(Z)^{\circ} \cap \mathrm{omb}(X')^{\circ}$. Prouvons que $X'$ est plus petit dans $\widetilde{X}^{Z}$, i.e.

\page{29}
\[
  Z \ll X'' \leq X \Longrightarrow X' \leq X'' .
\]
Soit donc \struck{$x \in \mathrm{omb}(Z)^{\circ}$ (ML 2)} comme $x \in \mathrm{omb}(X')^{\circ}$, \struck{$x \in \mathrm{omb}(X'')$, \ill{}} on a $\mathrm{omb}(X')^{\circ} \cap \mathrm{omb}(X'') \neq \emptyset$, d'où $X' \leq X''$ par lemme 1. OK.
\note{dans $\mathrm{omb}(X')^{\circ}$, le $X'$ est écrit en surcharge sur un $X$}

\emph{Corollaire 4} \quad $X \mathrel{\overset{\circ}{\ll}} Y \mathrel{\overset{\circ}{\ll}} Z \Longrightarrow X \mathrel{\overset{\circ}{\ll}} Z$
\note{le « 4 » paraît écrit en surcharge ; la page 32 renvoie à ce « cor.~4 p.~29 »}

Immédiat sur le critère (ii) du cor.~3
\marginal{At$_{\mathcal{M}}$ 5}

Il est temps maintenant de faire intervenir la relation $\parallel$ sur \struck{\ill{}} $\mathcal{L}$. On va \struck{\ill{}} \ill{} la condition minimal dans $x \in \mathcal{L}$ (?) (je pense) utiliser

\struck{At} Je rappelle
(ML 0) \quad la relation $x \parallel y$ \add{sur $\mathcal{L}$} est symétrique et antiréflexive (elle implique $x \neq y$)
\note{ce « ML 0 », pour la relation $\parallel$ sur $\mathcal{L}$, n'est pas le ML 0 de la page 25 ($X \leq Y \Rightarrow X \ll Y$). Dans la marge gauche, en oblique : « relation symétrique, \struck{réflexive} \ill{} », en partie biffé}

On pose dans $\mathcal{M}$
\[
  (1) \qquad X \between Y \overset{\mathrm{def}}{\iff} \forall\, x \in \mathrm{omb}(X) \smallsetminus \mathrm{omb}(L),\ y \in \mathrm{omb}(Y) \smallsetminus \mathrm{omb}(L) \ \text{on a}\ x \parallel y
\]
où
\[
  \mathrm{omb}(L) \overset{\mathrm{déf}}{=} \bigcup_{\substack{Z \in \widetilde{X} \cap \widetilde{Y} \\ \text{i.e.}\ Z \leq X, Y}} \mathrm{omb}(Z)
\]

\begin{tikzcd}[column sep=small, row sep=small]
\mathrm{omb}(X) \arrow[dr, no head] & & \mathrm{omb}(Y) \arrow[dl, no head] \\
 & \mathrm{omb}(X) \cap \mathrm{omb}(Y) \arrow[d, no head] & \\
 & \mathrm{omb}(L) &
\end{tikzcd}

\drawing{schéma d'inclusions tracé par simples traits ; un premier « $\mathrm{omb}(Y)$ » est biffé et récrit à côté}

Dans \ill{} \ill{}, par ML 0, \struck{$\mathrm{omb}(X), \mathrm{omb}(Y)$}
\[
  \mathrm{omb}(X) \cap \mathrm{omb}(Y) = \mathrm{omb}(L)
\]
\note{« ML 0 » est sur la page ; c'est ML 1 (page 25) qui donne cette égalité}

\emph{Lemme 3} Si $X, Y \leq Z$, on a $X \between Y$

Soit en effet \struck{$x \in \mathrm{omb}$}
\marginal{inutile ici, sera contenu dans lemme 5}
\note{note oblique dans la marge gauche, en regard du lemme 3 ; le lemme 5 est à la page 30}

\page{30}
\begin{tikzcd}[column sep=small, row sep=small]
 & Z & \\
X \arrow[ur, no head, "\leq" description] & & Y \arrow[ul, no head, "\leq" description] \\
 & L \arrow[ul, no head, "\leq" description] \arrow[ur, no head, "\leq" description] &
\end{tikzcd}

\note{losange tracé en tête de page (à gauche de son numéro « 30 ») : $X, Y \leq Z$, $L \leq X, Y$}

En effet, soit $x \in \mathrm{omb}(X) \smallsetminus \mathrm{omb}(L)$, je dis qu'on a $x \notin \mathrm{omb}(Y)$ --- en effet $\mathrm{omb}(X) \cap \mathrm{omb}(Y) = \mathrm{omb}(L)$ par ML 1 appliqué à $X, Y$. On utilisera \add{alors}

\emph{ML 4} \quad Soit $Y \leq X$ (\struck{\ill{}} ald $Y \leq Z$), et $x \in \mathrm{omb}(X) \smallsetminus \mathrm{omb}(Y)$ (comme $Y \neq X$ \ldots), $y \in \mathrm{omb}(Y)$, alors $x \parallel y$
\marginal{At L 3}
\note{dans « $Y \leq X$ », « $\mathrm{omb}(X)$ », « $\mathrm{omb}(Y)$ » et « $Y \neq X$ », les lettres sont écrites en surcharge sur d'autres ; la parenthèse « (ald $Y \leq Z$) », au-dessus de la ligne, est reliée par une flèche au $X$ de « $Y \leq X$ »}

\emph{Lemme 4} Soient $X, Y \in \mathcal{M}$, $X \between Y$. Alors l'inverse du lemme 1 est valable \struck{pour} $X, Y$. Il suffit encore de prouver
\[
  \mathrm{omb}(X)^{\circ} \cap \mathrm{omb}(Y) \neq \emptyset \Longrightarrow X \leq Y .
\]
\marginal{(N'utilise \ill{} ML 4 \ill{} At$_{\mathcal{M}}$ 6)}
\note{note oblique dans la marge gauche, en regard du lemme 4}

\struck{Prouvons} Soit donc $x \in \mathrm{omb}(X)^{\circ} \cap \mathrm{omb}(Y)$ i.e.
\[
  x \mathrel{\overset{\circ}{\ll}} X, \quad x \ll Y
\]
\struck{On ne peut avoir ici \ill{} $x \notin$ $x \in \mathrm{omb}(X)$}
En fait, comme $\mathrm{omb}(X) \cap \mathrm{omb}(Y) = \mathrm{omb}(L)$, $\exists\, Z$ avec
\[
  x \ll Z \leq X, Y
\]
comme $x \mathrel{\overset{\circ}{\ll}} X$, on a $Z = X$, donc $X \leq Y$, cqfd.

\emph{Lemme 5} Soit $X, Y, X', Y'$ avec \struck{\ill{}}, $X' \leq X$, $Y' \leq Y$. Alors $X \between Y \Longrightarrow X' \between Y'$.

Soit en effet
\[
  \begin{cases}
    x \in \mathrm{omb}(X') \smallsetminus \mathrm{omb}(L') \\
    y \in \mathrm{omb}(Y') \smallsetminus \mathrm{omb}(L') \\
    \qquad \text{où}\ \mathrm{omb}(L') \overset{\mathrm{déf}}{=} \bigcup_{Z \in \widetilde{X}' \cap \widetilde{Y}'} \mathrm{omb}(Z)
  \end{cases}
\]
prouvons $x \parallel y$.

\page{31}
Si $x, y \notin \mathrm{omb}(L)$, cela résulte de l'hypothèse $X \between Y$. Si \ill{} $x$ (p.\ ex.) est p.\ ex.\ $x \in \mathrm{omb}(L)$, \struck{\ill{}} i.e.
\[
  \exists\, Z \in \widetilde{X} \cap \widetilde{Y} \quad (\text{i.e.}\ Z \leq X, Y) \ \text{avec}\ x \ll Z
\]
[Appliquons ML 1 à

\begin{tikzcd}[column sep=small, row sep=small]
X' \arrow[r, no head, "\leq" description] & X \\
x \arrow[u, no head, "\ll" description] \arrow[r, no head, "\ll" description] & Z \arrow[u, no head, "\leq" description]
\end{tikzcd}

, on trouve \struck{$S \in \mathcal{M}$, $S \in \widetilde{X}$ \ill{}}

\begin{tikzcd}[column sep=small, row sep=small]
X' \arrow[r, no head, "\leq" description] & X & \\
S \arrow[u, no head, "\leq" description] \arrow[r, no head, "\leq" description] & Z \arrow[u, no head, "\leq" description] \arrow[r, no head, "\leq" description] & Y \\
x \arrow[u, no head, "\ll" description] & &
\end{tikzcd}

\note{sur la page, « $(\leq Y)$ » est écrit après $Z$ ; à droite, le même contenu est redessiné en « i.e. » : $x \ll S \leq X'$, $S \leq Y$}
\struck{donc $S \in \widetilde{L} \overset{\mathrm{déf}}{=} \widetilde{X} \cap \widetilde{Y}$]}

\struck{donc $x \ll$ \ill{}} $x \ll Y$, \add{i.e.\ $x \in \mathrm{omb}(Y)$}, et pour prouver \struck{Je dis} $x \parallel y$, il suffit par \emph{ML 4} de prouver $x \notin \mathrm{omb}(Y')$. Mais si on avait $x \in \mathrm{omb}(Y')$, on aurait \struck{i.e.\ $x \ll$} par ML 1 appliqué à

\begin{tikzcd}[column sep=small, row sep=small]
S \arrow[r, no head, "\leq" description] & Y \\
x \arrow[u, no head, "\ll" description] \arrow[r, no head, "\ll" description] & Y' \arrow[u, no head, "\leq" description]
\end{tikzcd}

un $T$ tel que
\[
  x \ll T \leq S, Y'
\]
et on aurait donc aussi
\[
  T \leq X', Y' \qquad \text{puisque}\ S \leq X'
\]
donc $x \in \mathrm{omb}(L')$, contrairement à l'hyp., cqfd.

\struck{\ill{}} Il faut prouver maintenant \struck{\ill{}} seulement At$_{\mathcal{M}}$ 4 \struck{At$_{\mathcal{M}}$ \ill{}} b) (At$_{\mathcal{M}}$ 5, 6 et At$_{\mathcal{M}}$ 4 a) sont déjà prouvés) pour conclure que $(\mathcal{M}, \leq, \ll, \underset{\mathcal{M}}{\between})$ (et \struck{la famille} $\mathfrak{F}$ = ens.\ des parties fermées de $\mathcal{M}$ formées d'él.\ 2 à 2 compatibles pour $\underset{\mathcal{M}}{\between}$)

\page{32}
\struck{\ill{}} correspondent à un atelier At 1 -- At 7. Donc c'est le contenu du

\emph{Lemme 6} Soient $X, Y, Z$ avec
\[
  Z \mathrel{\overset{\circ}{\ll}} X, \quad Z \mathrel{\overset{\circ}{\ll}} Y, \quad X \between Y
\]
alors $X = Y$.

Ça résultera du cor.\ plus général

\emph{Cor} On a, si $X \between Y$
\[
  Z \mathrel{\overset{\circ}{\ll}} X, \ Z \ll Y \Longrightarrow X \leq Y
\]
\note{le corollaire est marqué d'un trait vertical dans la marge ; le $\leq$ final porte une petite marque, peut-être le rond de $\overset{\circ}{\ll}$}

\emph{Dém} Soit $x \in \mathrm{omb}(Z)^{\circ}$ (ML 2). On a donc \struck{$x \mathrel{\overset{\circ}{\ll}}$}
\[
  x \mathrel{\overset{\circ}{\ll}} Z \mathrel{\overset{\circ}{\ll}} X
\]
donc (lemme 2, cor.~4 p.~29) $x \mathrel{\overset{\circ}{\ll}} X$, donc $x \in \mathrm{omb}(X)^{\circ} \cap \mathrm{omb}(Y)$, lequel est $\neq \emptyset$, donc $X \leq Y$ (lemme 4).

Il faut voir que la relation de départ $\parallel$ sur $\mathcal{L}$ est bien la relation de disjonction au sens de l'atelier qu'on vient de construire, i.e.

\emph{Lemme 7} \struck{\ill{}} a) Soient $x, y \in \mathcal{L}$. Alors
\[
  x \between y \ \text{ssi}\ x = y \ \text{ou}\ x \parallel y
\]
b) On a $x \parallel y$ ssi $x \between y$ et $\mathrm{Inf}^{\leq}(x, y) = \emptyset_{\mathfrak{F}}$. (i.e.\ $x$ et $y$ disjoints au sens de l'atelier)

Dans a), il suffit de voir le cas $x \neq y$. À cause du \emph{caractère minimal} de $x, y$, on a $\mathrm{omb}\, x = \lbrace x \rbrace$, $\mathrm{omb}(y) = \lbrace y \rbrace$, $L = \emptyset$, et la condition \struck{est bien} qui définit $x \between y$ est bien $x \parallel y$. L'assertion b) résulte,
\note{« $\mathrm{omb}\, x = \lbrace x \rbrace$, $\mathrm{omb}(y) = \lbrace y \rbrace$ » est écrit dans la marge gauche et rattaché à la ligne}

\page{33}
\emph{Lemme 8} Soient $X, Y \in \mathcal{M}$. Alors
\[
  \underbrace{X \between Y \ \text{et}\ \widetilde{X} \cap \widetilde{Y} = \emptyset}_{\text{i.e.}\ X \parallel Y \ \text{au sens de l'atelier}}
  \iff \forall\, x \in \mathrm{omb}(X),\ y \in \mathrm{omb}(Y) \Longrightarrow x \parallel y
\]
\marginal{At L 2}

C'est tautologique !

\struck{Lemme 9} Ainsi, notre atelier satisfait, en plus de At 1 -- At 7, \struck{les} conditions sur les lieux At L 1, At L 2, At L 3. Reste à voir la condition At CL (spéc) --- mais elle est vraie par définition.

On trouve ainsi le

\emph{Théorème-Scholie} Soit $(\mathcal{M}, \leq, \ll)$ un ensemble muni de deux relations d'ordre, $\mathcal{L}$ l'ens.\ des éléments minimaux de $\mathcal{M}$ pour $\ll$, $\parallel$ une relation symétrique antiréflexive sur $\mathcal{L}$. Pour que \add{ces données} \struck{proviennent} soient isom.\ à celles provenant d'un atelier \add{spécieux} satisfaisant \struck{At 1} (en plus de At 1 -- At 7) les axiomes sur les lieux At L 1, At L 2, At L 3, \struck{plus} At CL (spéc), il faut et suffit que les conditions suivantes soient satisfaites :

\page{34}
(ML 0) \quad $X \leq Y \Longrightarrow X \ll Y$

(ML 1) \quad Soient $Z \ll X \geq X'$ dans $\mathcal{M}$, et soit $x \in \mathcal{L}$ avec $x \ll Z$, $x \ll X'$. Alors $\exists\, Z' \in \mathcal{M}$ avec

\begin{tikzcd}[column sep=small, row sep=small]
 & Z \arrow[r, no head, "\ll" description] & X \\
x \arrow[r, no head, "\ll" description] & Z' \arrow[u, no head, "\leq" description] \arrow[r, no head, "\ll" description] & X' \arrow[u, no head, "\leq" description]
\end{tikzcd}

De plus, si $Z \leq X$, on peut supposer $Z' \leq X'$.
\marginal{axiome trivial}
\marginal{axiome des raffinements \uncertain{induits}}
\note{les deux notes sont écrites en oblique dans la marge gauche, en regard de ML 0 et de ML 1 ; les deux « (ML 0) », « (ML 1) » sont écrits sur un premier sigle biffé. Sur la page, $X'$ est sous $X$, relié par un $\geq$ vertical}

(ML 2)* \quad $\forall\, X \in \mathcal{M}$, $\mathrm{omb}(X)^{\circ} \neq \emptyset$ i.e.\ $\exists\, x \in \mathcal{L}$, avec $x \mathrel{\overset{\circ}{\ll}} X$ i.e.\ $x \ll X$ et $(x \ll X' \leq X) \Rightarrow X' = X$

(ML 3) \quad \struck{\ill{}} $\forall\, X \in \mathcal{M}$, $x \in \mathcal{L}$ avec $x \ll X$, $\exists\, X' \leq X$ avec $x \mathrel{\overset{\circ}{\ll}} X'$
\marginal{axiome de divisibilité \struck{\ill{}} \ill{} \ill{} \ill{} et \struck{\ill{}} \ill{} la plus petite strate}
\note{note oblique dans la marge gauche, en regard de ML 2 -- ML 3, en partie biffée ; « la plus petite strate » est lu avec doute}

(ML 4) \quad $\parallel$ est symétrique antiréflexive ($x \parallel y \Rightarrow y \parallel x$ et $x \neq y$)

(ML 5) \quad $\forall\, X, Y \in \mathcal{M}$ avec $Y \leq X$, \struck{\ill{}} $y \in \mathrm{omb}(Y)$, $x \in \mathrm{omb}(X) \smallsetminus \mathrm{omb}(Y) \Longrightarrow x \parallel y$
\marginal{axiomes de disjonction pour les lieux}
\marginal{c'est tout, des axiomes \ill{} \ill{} sur la relation $\parallel$}
\note{la première note est à droite de ML 4, avec une accolade ; la seconde, oblique, est dans la marge gauche, avec une accolade réunissant ML 4 et ML 5. Le ML 4 de cette page est le ML 0 de la page 29, et son ML 5 le ML 4 de la page 30 : il renumérote}

\struck{Les ateliers correspondants à ces \uncertain{types} de données, réalisés comme des ens.\ de parties de $\mathcal{M}$, $\mathfrak{F} \subset \mathfrak{P}(\mathcal{M})$, \ill{} les $\mathfrak{F}$ satisfaisant les conditions suivantes (i)}
\note{ce paragraphe inachevé est barré de six longs traits obliques}

\note{l'astérisque de ML 2 renvoie à la note qui suit, en bas de page}
* \emph{NB} Cet axiome implique que si les relations d'ordre $\leq$ et $\ll$ sont égales, alors $\mathcal{M}$ est discret i.e.\ $\mathcal{M} = \mathcal{L}$. \struck{\ill{} \ill{} \ill{}} L'axiome ML 4 est vide !
\note{« ML 4 » est sur la page ; dans la numérotation de cette page, c'est ML 5 qui est vide quand $\leq$ est discret ($Y \leq X$ force $Y = X$), ML 4 ne l'étant pas : il s'agit sans doute de l'ancien ML 4 de la page 30}

\page{35}
Soit \struck{\ill{}} $\underline{\Phi}$ l'ens.\ des parties de $\mathcal{M}$ qui sont a) fermées pour $\leq$ et b) formées d'éléments deux à deux compatibles (au sens de (1) p.~29). Munissons $\underline{\Phi}$ de la relation $\leq \overset{\mathrm{déf}}{=}$ l'inclusion, et $\ll$ défini par
\[
  F \ll G \overset{\mathrm{déf}}{\iff} \forall\, X \in F,\ \exists\, Y \in G \ \text{avec}\ X \ll Y .
\]
\note{$\underline{\Phi}$ rend une capitale en forme de $\Phi$ soulignée, écrite après un mot raturé surmonté d'un astérisque ; même sigle dans toute la suite de la page}
Alors $\underline{\Phi}$ est un atelier \struck{\ill{}} satisfaisant \struck{At} At L 1, At L 2, At L 3 et At CL (spéc), et \ill{} les $\widetilde{X} = \mathcal{M}_{\leq X}$ \struck{\ill{}} ($X \in \mathcal{M}$) comme \struck{objets} figures élémentaires, \struck{les $\lbrace x \rbrace$ \ill{} et} \ill{} induisant sur $\mathcal{M}$ les relations \add{(\struck{\ill{}} de disjonction des figures)} \struck{\ill{}} \add{et} la relation $\parallel$ \struck{\ill{}} $\leq$, $\ll$ de départ, et \ill{}

\struck{Les autres} \add{Tout atelier} correspondant à ces \struck{\ill{}} données s'identifie : \emph{aux sous-ateliers \uncertain{spécieux}} de $\underline{\Phi}$, \struck{i.e.\ parties $\mathfrak{F}$} \struck{de $\underline{\Phi}$} i.e.\ une partie
\[
  \mathfrak{F} \subset \underline{\Phi}
\]
satisfaisant
\begin{itemize}
  \item[(i)] $\forall\, F \in \mathfrak{F}$, et $G \leq F$ \struck{\ill{}} (i.e.\ $G$ partie fermée de $F$, $\in$) on a $G \in \mathfrak{F}$ et
  \item[(ii)] $\forall\, F, G \in \mathfrak{F}$, si $F \between G$ i.e.\ $F \cup G \in \underline{\Phi}$ (i.e.\ \struck{\ill{}} éléments de $F$ sont compatibles à tous \uncertain{éléments} de $G$) $F \cup G \in \mathfrak{F}$.
\end{itemize}
\note{« sous-ateliers » et le mot qui suit sont soulignés ; ce dernier est lu « spécieux » avec doute}

\page{36}
Il contient de plus les \struck{objets} figures élémentaires, i.e.\ satisfait :
\begin{itemize}
  \item[(iii)] $\forall\, X \in \mathcal{M}$, $\lbrace\lbrace X \rbrace\rbrace \in \mathfrak{F}$
\end{itemize}
\note{ainsi sur la page, avec des accolades doublées ; d'après la page 35, la figure élémentaire attendue est $\widetilde{X}$}
et par suite, il contient toutes les parties fermées de réunions finies de tels $\widetilde{X}$.

Inversement, tout $\mathfrak{F}$ satisfaisant ces trois conditions, satisfait aux conditions voulues \struck{\ill{}} pour que les figures de $\mathfrak{F}$ soient \ill{}, \ill{} et soient \struck{soit}
\[
  a)\ \mathfrak{F} = \mathfrak{F}_0 = \lbrace \text{réunions finies d'ensembles de la forme}\ \widetilde{X} \rbrace
\]
(ce qui \ill{} donne un atelier qui \ill{} supposé de plus
\[
  b)\ \text{toute partie fermée d'un}\ \widetilde{X}\ \text{\add{(pour $\leq$)} est de type fini}
\]
À part donc un petit flottement dû à la possibilité de \struck{familles \ill{} infinies} \ill{} infinies, on voit que \struck{\ill{}} \add{tout} s'exprime par $(\mathcal{M}, \leq, \ll, \parallel)$, avec $\parallel$ défini sur $\mathcal{L}$ sans plus, \struck{\ill{} \ill{}} satisfaisant aux quatre axiomes (cinq avec (ML 0)) ML 1 -- ML 4 (au lieu
\note{« quatre axiomes \ldots\ ML 1 -- ML 4 » : la page 34 en énumère six, ML 0 -- ML 5 ; la phrase se poursuit à la page 37}

\page{37}
de la douzaine d'axiomes At 1 -- 7, At L 1 -- 3, At CL (spéc).

Si on se donne $\parallel$ sur $\mathcal{M}$, et non sur $\mathcal{L}$, il faut simplement rajouter un axiome supplémentaire :
\[
  (\mathrm{ML}\ 6) \quad \text{Soit}\ X, Y \in \mathcal{M}.\ \text{Alors} \quad X \parallel Y \iff \forall\, x, y \in \mathcal{L},\ x \ll X,\ y \ll Y \Rightarrow x \parallel y .
\]
\note{le « 6 » de ML 6 est repassé}

\emph{Variante pour ateliers polyédraux.} (\add{satisfaisant At 1 -- At 7, At L 1 -- At L 3, (At CL (pol))}) Elle consiste en le \ill{} texte, sauf qu'on change un \uncertain{peu} la définition de $X \between Y$ (p.~29), et qu'il faut supposer sur $\mathcal{M}$, $\leq$ une condition supplémentaire
\note{l'addition entre parenthèses est écrite au-dessus de la ligne, dans un cadre arrondi relié à « polyédraux »}

(ML 0 pol) \quad $\forall\, X \in \mathcal{M}$, $\widetilde{X} = \mathcal{M}_{\leq X}$ est un ens.\ ordonné \emph{polyédral} (p.~13), i.e.\ \struck{$\exists\, X, Y, Z$} deux élts $Y, Z$ de \struck{$\mathcal{M}$} \add{$\widetilde{X}$ \uncertain{qui}} sont majorés (par un $X$) et minorés \struck{\ill{}} (par un $T$) admettent un inf \add{$\leq$} --- en d'autres termes, les ensembles $\mathcal{M}_{T \leq \cdot \leq X}$ ont des inf finis (donc aussi des sup finis s'ils sont finis).
\note{le paragraphe (ML 0 pol) est marqué d'un trait vertical dans la marge ; la page 13 est au lot 1}

\page{38}
Ça montre, de plus, que si on part de $(\mathcal{M}, \leq, \ll)$, il y a un très large choix pour les $\parallel$ possibles sur $\mathcal{L}$, i.e.\ les parties \uncertain{admissibles} de $\mathfrak{P}_2(\mathcal{L})$. Si on désigne par
\[
  \Gamma_0 = \Bigl\lbrace \lbrace x, y \rbrace \in \mathfrak{P}_2(\mathcal{L}) \Bigm| \exists\, X, Y \in \mathcal{M},\ Y \leq X,\ y \ll Y,\ x \ll X,\ x \not\ll Y \Bigr\rbrace
\]
\note{dans « $y \ll Y$ », le $Y$ est écrit en surcharge}
on peut prendre comme structure de graphe $\Gamma$ \struck{de} (pour une relation $\parallel$) n'importe quelle partie $\Gamma \supset \Gamma_0$.
\[
  \Gamma_0 \subset \Gamma \subset \mathfrak{P}_2(\mathcal{L}) .
\]
Ça fait de la marge ! Pour la suite, il faudra voir comment les \uncertain{traduire} \ill{} les \uncertain{choix} de \uncertain{certains} axiomes \ill{} \ill{} le choix de $\parallel$ (et des variantes). Pour le moment, je songe surtout :
\note{les mots « traduire », « choix » et « certains » de cette phrase sont lus avec doute}

At supp \quad cf p.~18)

At ens et ses variantes \quad (cf GF VI, p.~52 \ldots)
\note{la page 18 est au lot 1. « GF VI, p.~52 » renvoie à la deuxième mouture (dossier 156-6), où la page 52 définit l'atelier ensembliste et l'axiome « At ens 1 » ; le « VI » est souligné}

\struck{Donc} Pour qu'un atelier \add{spécieux} \add{(resp.\ polyédral)} ensembliste, \add{suffit-il} que $x \parallel y \Leftrightarrow x \neq y$ ?? Sans doute pas, car cela n'implique aucune restriction sur les autres données $\leq$, $\ll$, et les conditions ML 0 à ML 3 n'impliquent pas que l'application voulue : $\mathcal{M} \longrightarrow \mathrm{Fig}(\mathcal{L})$ soit injective, ni que $\ll$ soit induite par la relation $\ll$ entre fig.\ ensemblistes dans $\mathcal{L}$.
\note{ce dernier alinéa est marqué d'un double trait vertical dans la marge ; « (resp.\ polyédral) » et « suffit-il » sont écrits au-dessus et au-dessous de « Pour qu'un », dans un cadre}

\page{39}
Dans le raisonnement précédent, le fait que les éléments de $\mathcal{L}$ soient minimaux n'est intervenu qu'à la fin, pour prouver que la relation de départ $\parallel$ sur $\mathcal{L}$ est bien celle induite par la disjonction des lieux, au sens de l'atelier construit. On va tout reprendre.

\emph{Définition} \struck{Soit} $\Lambda \subset \mathcal{M}$. On dit que $\Lambda$ est \emph{dense} \add{(dans $\mathcal{M}$)} si $\forall\, X$ \struck{\ill{}} $\in \mathcal{M}$, $\exists\, Y \in \Lambda$ avec $Y \mathrel{\overset{\circ}{\ll}} X$.
\note{la définition est marquée d'un trait vertical dans la marge}

\emph{Ex} 1${}^{\circ}$) $\Lambda = \mathcal{M}$, exemple « universel » (on prend $Y = X$)

2${}^{\circ}$) $\Lambda = \mathcal{L}$ \struck{\ill{}} quand At L 1 est satisfait

On va donner des axiomes sur $(\mathcal{M}, \leq, \ll, \between, \Lambda \subset \mathcal{M})$ qui assurent que $(\mathcal{M}, \leq, \ll, \between)$ provient d'un atelier, et que $\Lambda \subset \mathcal{M}$ est dense, avec \struck{\ill{}} en plus sur l'atelier (au besoin) la condition At D ou At D$_0$ ou At C.

(M$\Lambda$ 0) \quad $X \leq Y \Longrightarrow X \ll Y$
\marginal{Axiome trivial}

(M$\Lambda$ 1) \quad Si \struck{$Z$} $Y \ll X \geq X'$ et \struck{\ill{}} $Z \in \Lambda$ \ldots, $Z \ll Y$, $Z \ll X'$, \struck{\ill{}} alors $\exists\, Y'$ avec \struck{$Z' \leq Z$,} \struck{$Y'$} $Y' \leq Y$, $Y' \ll X'$, $Z \ll Y'$ et si $Y \leq X$, $Y' \leq X'$.
\note{ligne très surchargée, lue avec doute : l'hypothèse porte sur $Y \ll X \geq X'$ (le $Y$ est écrit sur un $Z$ biffé) et sur $Z \in \Lambda$ (souligné), avec $Z \ll Y$, $Z \ll X'$ ; dans la conclusion, les $Y'$ sont écrits sur des lettres biffées. Une flèche en pointillé mène de la fin de la ligne de M$\Lambda$ 0 à une accolade dans la marge droite, où quelques mots obliques ne sont pas lus}
\marginal{Axiome du raffinement induit}

(M$\Lambda$ 2) \quad $\forall\, X \in \mathcal{M}$, $\exists\, Z \in \Lambda$ tel que $Z \mathrel{\overset{\circ}{\ll}} X$
\marginal{Axiome de densité \struck{\ill{}}}

(M$\Lambda$ 3) \quad $\forall\, X \in \mathcal{M}$, $Z \in \Lambda$ avec $Z \ll X$, $\exists\, X' \leq X$ avec $Z \mathrel{\overset{\circ}{\ll}} X'$
\marginal{Axiome de la \uncertain{strate} minimale}

\struck{Il} Il est commode d'introduire les
\[
  \begin{aligned}
    \mathrm{omb}_{\Lambda}(X) &= \lbrace Z \in \Lambda \mid Z \ll X \rbrace \\
    \mathrm{omb}_{\Lambda}(X)^{\circ} &= \lbrace Z \in \Lambda \mid Z \mathrel{\overset{\circ}{\ll}} X \rbrace .
  \end{aligned}
\]
Les raisonnements précédents montrent la validité des deux lemmes suivants, que je rappelle pour mémoire :
\marginal{Pour ce \ill{} il s'agit \ill{} axiomes \ill{} $\mathcal{L}$ \ill{} \ill{}}
\note{les quatre notes de marge sont écrites en oblique dans la marge gauche, en regard de M$\Lambda$ 0, M$\Lambda$ 1, M$\Lambda$ 2 et M$\Lambda$ 3 ; la dernière, plus bas, n'est lue qu'en partie}

\page{40}
\emph{Lemme A} Soient $X, Y \in \mathcal{M}$, supposons $\exists\, Z$ avec $X, Y \leq Z$. Conditions équivalentes
\begin{itemize}
  \item[(i)] $X \leq Y$
  \item[(ii)] $X \ll Y$
  \item[(ii bis)] $\mathrm{omb}_{\Lambda}(X) \subset \mathrm{omb}_{\Lambda}(Y)$
  \item[(iii)] $\mathrm{omb}_{\Lambda}(X)^{\circ} \subset \mathrm{omb}_{\Lambda}(Y)$
  \item[(iv)] $\mathrm{omb}_{\Lambda}(X^{\circ}) \cap \mathrm{omb}_{\Lambda}(Y) \neq \emptyset$
\end{itemize}
\marginal{cf lemme 1 p.~25}
\note{dans « $X, Y \leq Z$ », le $\leq$ est repassé ; dans (iv), le rond est écrit sur $X$ et non sur $\mathrm{omb}_{\Lambda}(X)$}

\emph{Lemme B} Soit $Y \ll X$, alors $\exists\, X'$ avec $Y \mathrel{\overset{\circ}{\ll}} X' \leq X$.
\marginal{cf cor.~1 p.~28}

\emph{Lemme C} Soit $Y \ll X \geq X'$. Conditions équival.
\begin{itemize}
  \item[(i)] $Y \ll X'$
  \item[(i bis)] $\mathrm{omb}_{\Lambda}(Y) \subset \mathrm{omb}_{\Lambda}(X')$
  \item[(ii)] $\mathrm{omb}_{\Lambda}(Y)^{\circ} \subset \mathrm{omb}_{\Lambda}(X')$
  \item[(iii)] $\mathrm{omb}_{\Lambda}(Y)^{\circ} \cap \mathrm{omb}_{\Lambda}(X') \neq \emptyset$
\end{itemize}
\marginal{cf cor.~2 p.~28}
\note{$X'$ est écrit sous $X$, relié par un $\geq$ vertical}

\emph{Corollaire} Mêmes notations. Cond.\ équivalentes
\begin{itemize}
  \item[(i)] $Y \mathrel{\overset{\circ}{\ll}} X'$
  \item[(ii)] $\mathrm{omb}_{\Lambda}(Y)^{\circ} \subset \mathrm{omb}_{\Lambda}(X')^{\circ}$
  \item[(iii)] $\mathrm{omb}_{\Lambda}(Y)^{\circ} \cap \mathrm{omb}_{\Lambda}(X')^{\circ} \neq \emptyset$
  \item[(iv)] $X'$ est minimal dans $\widetilde{X}^{Y} = \lbrace X' \in \widetilde{X} \mid Z \ll X' \rbrace$
  \item[(v)] $X'$ est plus petit élément de $\widetilde{X}^{Y}$.
\end{itemize}
\marginal{cf cor.~3 p.~28}
\note{les quatre renvois de la marge gauche vont aux pages 25 et 28 de ce lot. Dans (iv), le $Y$ de $\widetilde{X}^{Y}$ est écrit sur un autre signe, et la condition garde « $Z \ll X'$ » de la page 28}

Nous faisons maintenant intervenir $\between$, et allons voir quelles propriétés de $\between$ font marcher nos raisonnements précédents (comme \struck{\ill{}} \uncertain{naguère} : $\parallel$). On pose ici, pour $X, Y \in \mathcal{M}$
\[
  X \parallel Y \overset{\mathrm{déf}}{\iff} X \between Y, \ \text{et}\ \widetilde{X} \cap \widetilde{Y} = \emptyset
\]
\note{devant $\widetilde{X} \cap \widetilde{Y}$, un signe raturé}

\emph{M$\Lambda$ 0'} \quad $\between$ est réfl.\ et sym. Si $X \between Y$, et $X' \leq X$, $Y' \leq Y$, on a $X' \between Y'$. \add{avec \struck{\ill{}} $X', Y' \in \Lambda$}
\marginal{axiomes triviaux pour $\between$}
\note{l'addition finale est entourée ; le « 0' » est écrit sur un autre signe}

\struck{\emph{M$\Lambda$ 4} Soient $X' \ll X$, $Y' \ll Y$ (supposons $X'_{Y} = \emptyset$ \ldots), $\nexists\, X'' \leq X'$, avec $X'' \ll Y$. Alors $X'' \parallel Y'$.}
\marginal{Axiome de disjonction}
\note{l'énoncé M$\Lambda$ 4 est barré de deux longs traits obliques ; sur la page $Y$ est écrit sous $X$, relié par un signe $=$ ou $\parallel$ vertical, et plusieurs lettres sont en surcharge}

\end{document}
